\input{preamble}
\title{More than One Way to Skin a Cat:\\
Preserving Verified Modes in RLVR}
\author{Anonymous Authors}
\begin{document}
\maketitle

\begin{abstract}
When a mathematical agent reaches a dead end, a different route may offer a
way forward. Yet learning to produce correct answers can quietly narrow the
alternatives it has left to try. We introduce \mb{} to make this loss visible.
Across five executable domains, it gives verified outputs exact identities,
so we can measure how many distinct successes the model still produces.
ReplayMaxRL turns these identities into memory, giving each stored mode equal
replay weight even when fresh samples rarely produce it. Across model scales, this replay
mechanism improves held-out correctness and output support under both Dr.GRPO
and MaxRL, including a complete harder graph-coloring demonstration.
The results motivate learning to preserve successful solutions alongside
learning to find them.
\end{abstract}

\section{Preserving alternatives for mathematical agents}
\label{sec:introduction}\label{sec:collapse}
Consider a mathematical agent whose first attempt fails a check. It may repair
that attempt or try a different route. A proof idea that is hard to complete
may give way to another; a different computation may succeed under constraints
that defeated the first. Tools can help it explore these routes, while verifiers check the results
\citep{brown2024monkeys,snell2024testtime}. Yet a verifier cannot select an
alternative the model never proposes.

Correctness feedback can hide the loss of these options. Reinforcement learning
with verifiable rewards (RLVR) reinforces correct outputs \citep{guo2025deepseekr1}.
A familiar solution earns the same correctness reward as a newly discovered
one. Accuracy can therefore improve even as the range of verified outputs shrinks. In our runs, these two trends occur together
(Figure~\ref{fig:story}).

\begin{figure}[!htbp]
\centering
\includegraphics[width=\linewidth]{figures/modecollapse_story.pdf}
\caption{\textbf{Correctness can rise while alternatives disappear.}
One graph admits six valid colorings. Bars pool four eight-sample draws from
matched Qwen2.5-3B runs. At pass 6.5, Dr.GRPO produces 22/32 correct from one
mode; ReplayMaxRL produces 32/32 across four.}
\label{fig:story}
\end{figure}

The loss also appears across our datasets. Averaged over five domains, eight
passes of Dr.GRPO reduce the initial number of verified alternatives beyond
the first correct output in eight samples by $99.4\%$ at Qwen2.5-0.5B,
$71.6\%$ at Falcon3-1B, and $41.5\%$ at Qwen2.5-3B
(Appendix Figure~\ref{fig:baseline-precheck}).

Finding a successful output leaves a further question: will training keep
that success available? Our approach uses executable
identities to tell verified outputs apart, then reuses those identities to
remember alternatives after they become rare.

\section{\mb{} makes output identity executable}
\label{sec:modebench}
\mb{} assigns each accepted output an exact, executable identity: a canonical
key whose distinct values define verified \emph{modes}
(Figure~\ref{fig:modebench-examples}). Keys capture computation routes in
Countdown and MathIR, tested program behaviors in Python, and valid
configurations in Graph and PantryPlan. Cosmetic aliases merge; distinct routes
can share an answer (key rules: Appendix~\ref{app:domain-overview}). This measures
task-defined output support; it does not establish diversity of latent reasoning strategies.

\begin{figure}[!htbp]
\centering
\includegraphics[trim={0 9bp 0 0},clip,width=\linewidth]{figures/modebench_examples.pdf}
\caption{\textbf{One prompt, different verified outputs.}
Keys retain a color vector, normalized computation tree, tested return vector,
algebraic state trajectory, or ingredient support. Cosmetic aliases merge;
these keys identify task-defined outputs rather than free-form reasoning.}
\label{fig:modebench-examples}
\end{figure}

\section{ReplayMaxRL: learn from new samples and remember successes}
\label{sec:method}\label{sec:maxrl-method}\label{sec:replay}
These keys let us store distinct verified solutions and revisit them when
they are absent from fresh rollouts. ReplayMaxRL combines these replay
updates with MaxRL learning from new samples
(Figure~\ref{fig:verified-support-story}).

\begin{figure}[!htbp]
\centering
% Crop only blank PDF margins; keep the figure itself unchanged.
\includegraphics[trim={4bp 23bp 4bp 10bp},clip,width=\linewidth]{figures/verified_support_story.pdf}
\caption{\textbf{Discovery and preservation use the same verified samples.}
MaxRL learns from fresh rewards; replay retains exemplars by key and revisits
rare successes. Memory cannot protect a mode before it is discovered.}
\label{fig:verified-support-story}
\end{figure}

MaxRL trains policy $\pi_\theta$ for success within $k$ samples across sampling
budgets \citep{tajwar2026maxrl}. Each training prompt $x$ keeps one
policy-generated exemplar $e_c$ per admitted key in bank $\mathcal B_x$.
With $|e_c|$ its token length, recurrent replay minimizes
\begin{equation}
L_{\mathrm{replay}}(x)=-\frac{1}{|\mathcal B_x|}
\sum_{c\in\mathcal B_x}\frac{\log\pi_\theta(e_c\mid x)}{|e_c|}.
\label{eq:lm-verified-replay}
\end{equation}
Each retained key receives equal replay weight, independent of its frequency in
fresh rollouts. Replacing MaxRL with Dr.GRPO \citep{liu2025understanding},
which debiases GRPO \citep{shao2024deepseekmath}, yields ReplayDr.GRPO.
Appendix~\ref{app:algorithm} gives the full update;
Appendix~\ref{app:theory} analyzes collapse and retention in idealized categorical models.

\noindent\textbf{Related work.} Prior work studies diversity collapse
\citep{cui2025entropy,dang2025diversity}, sampling-aware objectives, and success
replay \citep{zhang2025rlep}. Our focus is execution-defined identity for
measuring and retaining verified alternatives. Appendix~\ref{sec:related}
gives the full discussion.

\section{Does memory preserve verified output support?}
\label{sec:experiments}\label{sec:results}
\texttt{pass@8} is the chance that eight samples contain any correct response;
\texttt{distinct@8} is their average number of distinct verified modes.
\emph{Extra modes}, \texttt{distinct@8}$-$\texttt{pass@8}, counts alternatives
beyond the first correct output. Both support measures remain coupled to correctness.

We register five seeds per comparison and train for eight passes. Each domain
has 128 held-out prompts, evaluated with four eight-sample draws and excluded
from the replay bank. Matched Dr.GRPO/MaxRL controls run the bank code with zero
replay gradient. Figure~\ref{fig:cross-scale-terminal-effects} also compares
reward normalization, on-policy diversity (UCPO; \citealp{lochab2026ucpo}),
success replay (RLEP-Dr), and semantic entropy. Full protocols are in
Appendix~\ref{app:experimental-design}.

\begin{figure}[!htbp]
\centering
\includegraphics[width=\linewidth]{figures/experiment1_retention_comparator_matrix.pdf}
\caption{\textbf{Replay improves correctness and verified support across scales.}
(A) ReplayDr.GRPO minus Dr.GRPO. (B) Qwen2.5-0.5B alternatives and the untrained
reference minus Dr.GRPO. Daggers mark four admissible pairs; others use five. Average uses shared seeds.
Left: probability differences. Right: key counts. Bold marks column maxima;
uncertainty: Appendix~\ref{app:supporting}.}
\label{fig:cross-scale-terminal-effects}
\end{figure}

\noindent\textbf{Replay improves held-out correctness and support.}
\label{sec:results-retention}
ReplayDr.GRPO improves both means in 14 complete model--domain blocks and one
four-seed block; \texttt{distinct@8} rises in all 74 admissible seed pairs
(one conflicting source excluded; Appendix~\ref{app:source-integrity}).
On Qwen2.5-0.5B Countdown, \texttt{pass@8} rises from $.48$ to $.67$ and
\texttt{distinct@8} from $.49$ to $1.64$. Across five Qwen2.5-0.5B domains, uniform versus
frequency-weighted replay adds $.318$ extra modes on average
(95\% paired bootstrap interval $[.272,.358]$); its correctness effect is
inconclusive (Appendix~\ref{app:frequency-replay-progress}).

\clearpage
\begin{figure}[!htbp]
\centering
\includegraphics[width=\linewidth]{figures/e118_all_scale_factorial_progress.pdf}
\caption{\textbf{Memory adds value beyond sampling-aware MaxRL.}
Each track connects the untrained model, fresh-objective training, and its replay
counterpart. Faint paths are seeds; bold paths average domains within seed.
Falcon Dr.GRPO uses four common admissible seeds; other tracks use five. The
Qwen2.5-3B row shows its completed Dr.GRPO/replay track.}
\label{fig:maxrl-factorial}
\end{figure}

\noindent\textbf{Replay improves both Dr.GRPO and MaxRL.}
\label{sec:results-maxrl}
Adding replay improves both cross-domain averages under either objective at
Qwen2.5-0.5B and Falcon3-1B (Figure~\ref{fig:maxrl-factorial}); Qwen2.5-3B
also benefits on its completed
Dr.GRPO track. For ReplayMaxRL, both endpoints improve in all five Qwen2.5-0.5B
domains; Falcon correctness improves in five and support in four.
At Qwen2.5-0.5B, ReplayMaxRL minus MaxRL is $.307$ \texttt{pass@8} and $.991$
\texttt{distinct@8}
(unadjusted 95\% paired $t$ intervals $[.208,.406]$, $[.815,1.167]$).
These cross-domain summaries are post-hoc; Falcon Dr.GRPO uses four common
seeds descriptively (Appendix~\ref{app:results-maxrl}).

\begin{figure}[!htbp]
\centering
\includegraphics[width=\linewidth]{figures/modebench_level_admission.pdf}
\caption{\textbf{Replay across difficulty levels with Qwen2.5-0.5B-Instruct.}
Open/filled markers: Level 1/2. (A) Frozen-model success.
(B) Interim means weight the five domains equally after averaging seeds.
Checkpoints match across levels and methods; Pantry has one seed at step 0
(Appendix~\ref{app:level2-interim}).}
\label{fig:level2-admission}
\end{figure}

\noindent\textbf{Replaying stored solutions also helps on harder problems.}
\label{sec:results-levels}
Level 2 preserves validators, keys, and valid-mode-count histograms while
increasing difficulty: Graph hides four vertices instead of three.
Only Graph has a complete five-seed, four-method block after eight passes.
Dr.GRPO/MaxRL reach $.20/.43$ \texttt{pass@8} and $.22/.55$
\texttt{distinct@8}; replay raises these to $.76/.74$ and $1.24/1.21$.
The five-domain comparison in Figure~\ref{fig:level2-admission}B remains interim
(Appendix~\ref{app:results-levels}).

\section{Conclusion and Limits}
\label{sec:conclusion}\label{sec:discussion}\label{sec:limitations}
\mb{} makes a missing part of evaluation measurable: the range of distinct,
verified outcomes a model can produce for one problem. This matters because
different approaches can have different strengths. One may succeed where
another fails, or remain useful when constraints change. Understanding diversity
therefore gives a fuller picture of a model's capabilities than accuracy alone.

Our results show that replay can preserve alternatives while improving
held-out correctness. These controlled tasks provide a starting point for
studying diversity in broader mathematical reasoning. Replay retains only discovered,
stored modes; individual-mode survival remains unmeasured
(Appendix~\ref{app:ablations}). We have not shown that broader support causes
accuracy gains or helps solve later problems. The next step is to test which
alternatives help on richer mathematical tasks, and when.

\label{main:lastpage}
\clearpage
\bibliographystyle{plainnat}
\bibliography{example_paper}
\clearpage
\appendix
\input{appendix}
\end{document}
